\documentclass[xcolor = {dvipsnames, table}]{beamer}

\definecolor{dotlightblue} {HTML}{ADD8E6}
\definecolor{dotlightcoral}{HTML}{F08080}

\RequirePackage[english]{babel}
\RequirePackage[utf8]{inputenc}

\RequirePackage{fancyvrb}

\RequirePackage{listings}
\lstset{
    language = C,
    basicstyle = \ttfamily,
}

% \set{}
\RequirePackage{braket}

\RequirePackage{amsmath}

\RequirePackage{tikz}

% for arrows inside listings
\usetikzlibrary{decorations.text,calc}
\newcommand{\tikzmark}[2]{%
     \tikz[overlay,remember picture]%
        \node[text=black, inner sep=2pt] (#1) {#2};%
}

\RequirePackage{algorithm}
\RequirePackage[noend]{algpseudocode}
\algnewcommand \algorithmicin{\textbf{in}}
\algnewcommand \In {\algorithmicin\ }
\algrenewtext{ForAll}[2]%
    {\algorithmicfor\ #1 \algorithmicin\ #2 \algorithmicdo}

% reduce the size of the font in algorithms slightly
\makeatletter
\renewcommand{\ALG@beginalgorithmic}{\small}
\makeatother

%Information to be included in the title page:
\title{Theory makes beautiful programs}
\author{Jørgen Kvalsvik}
\date{01/09/2022}

\begin{document}
\frame{\titlepage}

\begin{frame}
    Modified condition/decision coverage (MC/DC)

    \begin{itemize}
        \item Why care about coverage?
        \item How is it different from branch coverage?
        \item What does maths have to do with it?
    \end{itemize}
\end{frame}


\begin{frame}
    Why even care?
    \begin{columns}
        \begin{column}{0.5\textwidth}
            \includegraphics[width = \textwidth]{img/bureaucrat.png}
        \end{column}

        \begin{column}{0.5\textwidth}
            \begin{itemize}
                \item DO-178B/C (Level A)
                \item ISO26262 (ASIL D)
            \end{itemize}
        \end{column}
    \end{columns}
\end{frame}

\newcommand \rowhl {\rowcolor{Cyan!20}}

\begin{frame}[fragile]
    The problem with branch coverage
    \begin{columns}
        \begin{column}{0.5\textwidth}
            \begin{lstlisting}
if ((a && b) || c) {
    //
} else {
    //
}
            \end{lstlisting}
        \end{column}

        \begin{column}{0.5\textwidth}
            \begin{tabular}{c c c c}
                        a & b & c \\
                        \hline
                \rowhl  F & F & F & F \\
                \rowhl  F & F & T & T \\
                        F & T & F & F \\
                        F & T & T & T \\
                        T & F & F & F \\
                        T & F & T & T \\
                        T & T & F & T \\
                        T & T & T & T \\
            \end{tabular}
        \end{column}
    \end{columns}
\end{frame}

\begin{frame}
    Modified condition/decision coverage satisfied if:
    \begin{itemize}
        \item every entry and exit point has been invoked
        \item every basic condition has taken on all possible outcomes
        \item each basic condition has been shown to independently affect the
              decision’s outcome
    \end{itemize}
\end{frame}

\begin{frame}
    \begin{columns}
        \begin{column}{0.7\textwidth}
            \begin{itemize}
                \item every entry and exit point has been invoked
            \end{itemize}

            Branch coverage (or decision coverage)
        \end{column}

        \begin{column}{0.3\textwidth}
            \begin{tabular}{c c c c}
                        a & b & c \\
                        \hline
                \rowhl  F & F & F & F \\
                \rowhl  F & F & T & T \\
                        F & T & F & F \\
                        F & T & T & T \\
                        T & F & F & F \\
                        T & F & T & T \\
                        T & T & F & T \\
                        T & T & T & T \\
            \end{tabular}
        \end{column}
    \end{columns}
\end{frame}

\begin{frame}
    \begin{columns}
        \begin{column}{0.7\textwidth}
            \begin{itemize}
                \item every basic condition has taken on all possible outcomes
            \end{itemize}

            Condition coverage
        \end{column}

        \begin{column}{0.3\textwidth}
            \begin{tabular}{c c c c}
                        a & b & c \\
                        \hline
                \rowhl  F & F & F & F \\
                        F & F & T & T \\
                        F & T & F & F \\
                        F & T & T & T \\
                        T & F & F & F \\
                \rowhl  T & F & T & T \\
                \rowhl  T & T & F & T \\
                        T & T & T & T \\
            \end{tabular}
        \end{column}
    \end{columns}
\end{frame}

\begin{frame}
    \begin{columns}
        \begin{column}{0.7\textwidth}
            \begin{itemize}
                \item each basic condition has been shown to
                      \emph{independently} affect the decision’s outcome
            \end{itemize}
            Modified condition/decision coverage
        \end{column}

        \begin{column}{0.3\textwidth}
            \begin{tabular}{c c c c}
                        a & b & c \\
                        \hline
                \rowhl  F & F & F & F \\
                \rowhl  F & F & T & T \\
                        F & T & F & F \\
                        F & T & T & T \\
                \rowhl  T & F & F & F \\
                        T & F & T & T \\
                \rowhl  T & T & F & T \\
                        T & T & T & T \\
            \end{tabular}
        \end{column}
    \end{columns}
\end{frame}

\begin{frame}
    \begin{itemize}
        \item Testing all $2^N$ inputs would not give more information
        \item Requires $N+1$ test cases
        \item Can catch bad logic
    \end{itemize}

    % TODO: bad logic?
\end{frame}

\begin{frame}
    MC/DC is \emph{awful} without tooling
\end{frame}

\begin{frame}[fragile]
    I wrote a patch for gcc

    \begin{Verbatim}[fontsize = \footnotesize]
$ git log -n 1 --format=short --shortstat
Author: Jørgen Kvalsvik <jorgen.kvalsvik@woven-planet.global>

    Add condition coverage profiling

21 files changed, 2952 insertions(+), 27 deletions(-)
    \end{Verbatim}

    \begin{Verbatim}[fontsize=\footnotesize]
 gcc/tree-profile.cc |  +978
    \end{Verbatim}
\end{frame}

\begin{frame}[fragile]
    \begin{block}{Demo}
        \begin{lstlisting}[language = sh, basicstyle = \scriptsize\ttfamily]
$ gcc --coverage -fprofile-conditions
    demo.c -o demo
$ ./demo 0 0 0
$ ./demo 0 0 1
$ ./demo 1 0 0
$ gcov --conditions demo
$ cat demo.c.gcov

    if ((a && b) || c) {
condition outcomes covered 4/6
condition  0 not covered (true)
condition  1 not covered (true)
        \end{lstlisting}
    \end{block}

    \begin{block}{Question}
        Why is \lstinline{a = 1} not covered?
    \end{block}
\end{frame}

\begin{frame}
    \begin{itemize}
        \item each basic condition has been shown to \emph{independently}
              affect the decision’s outcome
    \end{itemize}

    \begin{block}{Definition}
        A condition independently affects the outcome if changing it while
        keeping the other values constant changes the outcome
    \end{block}
\end{frame}

\begin{frame}[fragile]
    \begin{columns}
        \begin{column}{0.5\textwidth}
            \begin{lstlisting}[basicstyle = \footnotesize\ttfamily]
    if ((a && b) || c) {
        //
    } else {
        //
    }
            \end{lstlisting}
        \end{column}

        \begin{column}{0.5\textwidth}
            \begin{tabular}{c c c c}
                        a & b & c \\
                        \hline
                 \rowhl F & F & F & F \\
                        F & F & T & T \\
                        F & T & F & F \\
                        F & T & T & T \\
                \rowhl  T & F & F & F \\
                        T & F & T & T \\
                        T & T & F & T \\
                        T & T & T & T \\
            \end{tabular}
        \end{column}
    \end{columns}

    \begin{block}{Observation}
        Changing \lstinline{a} does not change the decision.
    \end{block}
\end{frame}

\begin{frame}
    This effect is called \emph{masking}

    \begin{itemize}
        \item \lstinline{* || true}
        \item \lstinline{* && false}
    \end{itemize}
\end{frame}

\begin{frame}
    \begin{block}{Observation - commutation}
        Reversing a boolean expression does not change its truth table
        \begin{align*}
            (P \wedge Q) & \equiv (Q \wedge P) \\
            (P \vee Q)   & \equiv (Q \vee P)   \\
        \end{align*}
    \end{block}

    \begin{block}{Observation}
        Masked conditions are \emph{short circuited} in the reversed expression
    \end{block}
\end{frame}

\begin{frame}
    \centering
    \begin{columns}
        \begin{column}{0.3\textwidth}
            \centering
            \lstinline{(a && b) || c}
            \begin{tabular}{c c c c}
                a & b & c & \\
                \hline
                F & * & F & F \\
                F & * & T & T \\
                F & * & F & F \\
                F & * & T & T \\

                T & F & F & F \\
                T & F & T & T \\
                T & T & * & T \\
                T & T & * & T \\
            \end{tabular}
        \end{column}

        \begin{column}{0.3\textwidth}
            \centering
            \lstinline{c || (b && A)}
            \begin{tabular}{c c c c}
                c & b & a & \\
                \hline
                F & F & * & F \\
                F & F & * & F \\
                F & T & F & F \\
                F & T & T & T \\

                T & * & * & T \\
                T & * & * & T \\
                T & * & * & T \\
                T & * & * & T \\
            \end{tabular}
        \end{column}
    \end{columns}
\end{frame}

\begin{frame}
    \begin{columns}
        \begin{column}{0.3\textwidth}
            \centering
            \lstinline{(a && b) || c}
            \begin{tabular}{c c c c}
                a & b & c & \\
                \hline
                F & * & F & F \\
                F & * & T & T \\
                F & * & F & F \\
                F & * & T & T \\

                T & F & F & F \\
                T & F & T & T \\
                T & T & * & T \\
                T & T & * & T \\
            \end{tabular}
        \end{column}

        \begin{column}{0.3\textwidth}
            \centering
            \lstinline{c || (b && a)}
            \begin{tabular}{c c c c}
                c & b & a & \\
                \hline
                F & F & - & F \\
                T & - & - & T \\
                F & T & F & F \\
                T & - & - & T \\

                F & F & - & F \\
                T & - & - & T \\
                F & T & T & T \\
                T & - & - & T \\
            \end{tabular}
        \end{column}

        \begin{column}{0.3\textwidth}
            \centering
            \lstinline{(a && b) || c}
            \begin{tabular}{c c c c}
                a & b & c & \\
                \hline
                F & * & F & F \\
                - & - & T & T \\
                F & * & F & F \\
                - & - & T & T \\

                T & F & F & F \\
                - & - & T & T \\
                T & T & * & T \\
                T & T & * & T \\
            \end{tabular}
        \end{column}
    \end{columns}

    \begin{block}{Note}
        Row order in \lstinline{c || (b && a)} changed
    \end{block}
\end{frame}

\begin{frame}
    \centering\lstinline{(a && b) || c}
    \begin{columns}
        \begin{column}{0.3\textwidth}
            \centering
            \begin{tabular}{c c c c}
                        a & b & c & \\
                        \hline
                \rowhl  F & * & F & F \\
                \rowhl  - & - & T & T \\
                        F & * & F & F \\
                        - & - & T & T \\

                \rowhl  T & F & F & F \\
                        - & - & T & T \\
                \rowhl  T & T & * & T \\
                        T & T & * & T \\
            \end{tabular}
        \end{column}

        \begin{column}{0.3\textwidth}
            \centering
            \begin{tabular}{c c c c}
                        a & b & c \\
                        \hline
                \rowhl  F & F & F & F \\
                \rowhl  F & F & T & T \\
                        F & T & F & F \\
                        F & T & T & T \\
                \rowhl  T & F & F & F \\
                        T & F & T & T \\
                \rowhl  T & T & F & T \\
                        T & T & T & T \\
            \end{tabular}
        \end{column}
    \end{columns}
\end{frame}

\begin{frame}
    \begin{itemize}
        \item Start solving for simple examples - a general solution
            \emph{also} solves simple inputs, draw insights from these
        \item For every added (failing) case, understand \emph{why} it failed
            (i.e.  systematic properties)
        \item Theory-backed programs almost self-materialize, easy to write
        \item Russ Cox' Regular expressions can be simple and fast
    \end{itemize}
\end{frame}

\begin{frame}[fragile]
    \begin{columns}
        \begin{column}{0.4\textwidth}
            \begin{lstlisting}[basicstyle = \footnotesize\ttfamily]
if ((a && b) || c) {
    // t
} else {
    // f
}
// e
            \end{lstlisting}
        \end{column}
        \begin{column}{0.4\textwidth}
            \begin{lstlisting}[basicstyle = \footnotesize\ttfamily]
_a:
  if (a) goto _b
  else   goto _c
_b:
  if (b) goto _t
  else   goto _c
_c:
  if (c) goto _t
  else   goto _f
_t:
  goto _e
_f:
  goto _e
_e:
            \end{lstlisting}
        \end{column}
    \end{columns}
\end{frame}

\begin{frame}[fragile]
    \begin{columns}
        \begin{column}{0.4\textwidth}
            \begin{lstlisting}[basicstyle = \footnotesize\ttfamily]
if ((a && b) || c) {
    // t
} else {
    // f
}
// e
            \end{lstlisting}
        \end{column}
        \begin{column}{0.4\textwidth}
            \includegraphics[width = \linewidth, keepaspectratio]{graph/fig1.pdf}
        \end{column}
    \end{columns}
\end{frame}

\begin{frame}
    Control flow graph

    \begin{itemize}
        \item Directed graph
        \item Nodes are an uninterruptible sequence of instructions
        \item Edges are next possible paths of execution
        \item Edges are labelled fallthrough, true/false (conditional), complex
        \item Fallthrough and conditional are mutually exclusive
    \end{itemize}
\end{frame}

\begin{frame}
    First problem: finding decisions
\end{frame}

\begin{frame}[fragile]
    \begin{columns}
        \begin{column}{0.4\textwidth}
            \begin{lstlisting}[basicstyle = \footnotesize\ttfamily]
if ((a && b) || c) {
    // t
} else {
    // f
}
// e
            \end{lstlisting}
        \end{column}
        \begin{column}{0.4\textwidth}
            \includegraphics[width = \linewidth, keepaspectratio]{graph/fig2.pdf}
        \end{column}
    \end{columns}
\end{frame}

\begin{frame}
    \begin{equation*}
        \bigcup \set{Succ(v) \mid v \in B} = N[B]
    \end{equation*}
    \begin{equation*}
        N[B] = B \cup O_B
    \end{equation*}

    \begin{tabular}{l l}
        $B$    & is a decision (boolean expression) \\
        $O_B$  & is the \emph{outcome} of $B$ \\
        $N(B)$ & is the \emph{open neighborhood} of $B$ \\
        $N[B]$ & is the \emph{closed neighborhood} of $B$ \\
    \end{tabular}
\end{frame}

\begin{frame}[fragile]
    \begin{columns}
        \begin{column}{0.5\textwidth}
            \begin{lstlisting}[escapeinside={*@}{@*}, aboveskip = 5em]
if (*@\tikzmark{fst}{}@*(a && b) || c*@\tikzmark{lst}{}@*) {
*@\tikzmark{then}{}@*
} else {
*@\tikzmark{else}{}@*
}
            \end{lstlisting}

            \begin{tikzpicture}[overlay,remember picture]
                \node (fst) at ($(fst) + (0.0em, +1.75ex)$) {};
                \node (lst) at ($(lst) + (0.0em, +1.75ex)$) {};
                \node (top) at ([yshift = 3em]$(fst)!0.5!(lst)$)
                        {uninterruptible};
                \draw [->] (top) to (fst);
                \draw [->] (top) to (lst);

                \node (then) at ($(then) + (5.0em, 0.0em)$) {};
                \node (else) at ($(else) + (5.0em, 0.0em)$) {};
                \node (out) at ($(then)!0.5!(else) + (5.0em, 0.0em)$) {outcome};

                \draw [->] (out) to (then);
                \draw [->] (out) to (else);
            \end{tikzpicture}

            \only<1>{
                \begin{equation*}
                    \bigcup \set{Succ(v) \mid v \in B}
                \end{equation*}
            }

            \only<2>{
                \begin{equation*}
                    E(B) = \set{(u, v) \in E | u \in B, v \in N[B]}
                \end{equation*}
                All edges in $E(B)$ are conditional
            }

            \only<3>{
                \begin{equation*}
                    Succ(B_\Omega) = O_B
                \end{equation*}
            }
        \end{column}

        \begin{column}{0.4\textwidth}
            \only<1-2>{
                \includegraphics[
                    width  = \linewidth,
                    height = 0.8\textheight,
                    keepaspectratio,
                ]{graph/fig3.pdf}
            }
            \only<3>{
                \includegraphics[
                    width  = \linewidth,
                    height = 0.8\textheight,
                    keepaspectratio,
                ]{graph/fig4.pdf}
            }
        \end{column}
    \end{columns}
\end{frame}

\newcommand{\includegraph}[2]{%
    \includegraphics<#1>[
        width  = \linewidth,
        height = \textheight,
        keepaspectratio,
    ]{graph/#2.pdf}
}

\begin{frame}[fragile]
    \begin{columns}
        \begin{column}{0.4\textwidth}
            \begin{lstlisting}[basicstyle = \footnotesize\ttfamily]
if ((a && b) || c) {
    // t
} else {
    // f
}
            \end{lstlisting}
        \end{column}

        \begin{column}{0.4\textwidth}
            \includegraph{1}{fig5-1}
            \includegraph{2}{fig5-2}
            \includegraph{3}{fig5-3}
            \includegraph{4}{fig5-4}
            \includegraph{5}{fig5-5}
            \includegraph{6}{fig5-6}
            \includegraph{7}{fig5-7}
            \includegraph{8}{fig5-8}
            \includegraph{9}{fig5-9}
            \includegraph{10}{fig5-10}
            \includegraph{11}{fig5-11}
        \end{column}
    \end{columns}
\end{frame}

\begin{frame}[fragile]
    \begin{columns}
        \begin{column}{0.4\textwidth}
            \begin{lstlisting}[basicstyle = \footnotesize\ttfamily]
if (a && b) {
    if (c) {
        // p
    } else {
        // q
    }
} else {
    // r
}
            \end{lstlisting}
        \end{column}

        \begin{column}{0.5\textwidth}
            \includegraph{1}{fig6-1}
            \includegraph{2}{fig6-2}
            \includegraph{3}{fig6-3}
            \includegraph{4}{fig6-4}
            \includegraph{5}{fig6-5}
        \end{column}
    \end{columns}
\end{frame}

\begin{frame}[fragile]
    \begin{columns}
        \begin{column}{0.6\textwidth}
            \begin{lstlisting}[basicstyle = \footnotesize\ttfamily]
if (a && b) {
    if (c) {
        // p
    } else {
        // q
    }
} else {
    // r
}
            \end{lstlisting}
            \begin{equation*}
                C = (\textcolor{dotlightblue}{G}, \textcolor{dotlightcoral}{G'})
            \end{equation*}
            \begin{equation*}
                \begin{split}
                    \forall_e \in E(G) & \bullet cond(e) \\
                    \Rightarrow & O_B \subset N[G] \\
                    \Rightarrow & B \subseteq G
                \end{split}
            \end{equation*}
        \end{column}

        \begin{column}{0.4\textwidth}
            \includegraphics[
                width   = \linewidth,
                height  = \textheight,
                keepaspectratio,
            ]{graph/fig7.pdf}
        \end{column}
    \end{columns}
\end{frame}

\begin{frame}[fragile]
    \begin{columns}
        \begin{column}{0.6\textwidth}
            \begin{lstlisting}[
                basicstyle = \footnotesize\ttfamily,
                escapeinside={*@}{@*},
                aboveskip = 5em,
            ]
if (*@\tikzmark{a}{}@*a && *@\tikzmark{b}{}@*b) {*@\tikzmark{open}{}@*
    if (c) {*@\tikzmark{c}{}@*
        // p*@\tikzmark{p}{}@*
    } else {*@\tikzmark{right}{}@*
        // q*@\tikzmark{q}{}@*
    }
} else {
    // r*@\tikzmark{r}{}@*
}
*@\tikzmark{e}{}@*
            \end{lstlisting}

            \begin{tikzpicture}[overlay,remember picture]
                \coordinate (top) at ($(b) + (2.0em, 2.0em)$) {};
                \coordinate (right) at ($(right) + (+4.0em, 0.0em)$) {};
                \node (a) at (a.{north east}) {};
                \draw [->, blue]
                    (a)     to [bend left = 30]
                    (top)   to [bend left = 40]
                    (right) to [bend left = 45]
                    (r)
                ;

                \draw [->, blue]
                    (a)                         to [bend left = 30]
                    ($(top)  + (0em, 1em)$)     to [bend left = 60]
                    ($(open) + (2em, 0em)$)     to [bend left = 40]
                    (c.{north east})
                ;

                \node (b) at (b.{north east}) {};
                \coordinate (top) at ($(top) + (2.0em, 0.0em)$) {};
                \coordinate (right) at ($(right) + (+2.0em, 0.0em)$) {};
                \draw [->, purple]
                    (b)     to [bend left = 30]
                    (top)   to [bend left = 40]
                    (right) to [bend left = 60]
                    (r.{south})
                ;

                \draw [->, purple]
                    (b)                         to [bend left = 40]
                    ($(top)  + (0em, 1em)$)     to [bend left = 60]
                    ($(open) + (3em, 0em)$)     to [bend left = 40]
                    (c.{east})
                ;

            \end{tikzpicture}
            No path from \textbf{then} to \textbf{else}
        \end{column}

        \begin{column}{0.4\textwidth}
            \includegraphics[
                width  = \linewidth,
                height = \textheight,
                keepaspectratio,
            ]{graph/fig8.pdf}
        \end{column}
    \end{columns}
\end{frame}

\begin{frame}[fragile]
    \begin{columns}
        \begin{column}{0.6\textwidth}
            \begin{lstlisting}[
                basicstyle = \footnotesize\ttfamily,
                escapeinside={*@}{@*},
                aboveskip = 1em,
            ]
if (a && b) {*@\tikzmark{top}{}@*
    if (c) {
        // p
    } else {
        // q
    }
} else {*@\tikzmark{bot}{}@*
    // r
}
            \end{lstlisting}

            \visible<2->{
                \begin{tikzpicture}[overlay,remember picture]
                    \node (top) at (top) {};
                    \node (bot) at (bot) {};
                    \node (mid) at ($(top)!0.5!(bot) + (+5.0em, 0.0em)$)
                        {outcome};
                    \draw [->] (mid.{north west}) to (top);
                    \draw [->] (mid.{south west}) to (bot);
                \end{tikzpicture}
            }

            \alt<1-2>{
                \begin{equation*}
                    \begin{split}
                        \forall_v \in then(B) \bullet B \subset A(v) \\
                        \forall_v \in else(B) \bullet B \subset A(v) \\
                    \end{split}
                \end{equation*}
                where $A(v)$ are the ancestors of $v$
            }{
                \begin{equation*}
                    \begin{split}
                        B &= A_G(v_T) = A_G(v_F) \\
                          &= \bigcap \set { A_G(v) \mid v \in N(G)} \\
                    \end{split}
                \end{equation*}
                where
                \begin{tabular}{l}
                    $A_G(v) = A(v) \cap G$ \\
                    $\set{ v_T, v_F } = O_B$ \\
                \end{tabular}
            }
        \end{column}

        \begin{column}{0.4\textwidth}
            \includegraphics[
                width  = \linewidth,
                height = \textheight,
                keepaspectratio,
            ]{graph/fig8.pdf}
        \end{column}
    \end{columns}
\end{frame}

\begin{frame}[fragile]
    \begin{columns}
        \begin{column}{0.5\textwidth}
            \begin{lstlisting}[
                basicstyle = \footnotesize\ttfamily,
            ]
if (a && b) {
    if (c) {
        // p
    } else {
        // q
    }
} else {
    // r
}
            \end{lstlisting}

            \setbeamercovered{invisible}

            % https://tex.stackexchange.com/questions/142672/uncovering-lines-in-an-equation-in-split-environment
            \newcommand{\disponslide}[2]{%
                \alt<#1>{#2}{\phantom{#2}}
            }
            \begin{align*}
                    \onslide<2->
                    \disponslide{2-}{A(q)}  & \disponslide{2-}{= \set{c, b, a}} \\
                    \disponslide{3-}{A(p)}  & \disponslide{3-}{= \set{c, b, a}} \\
                    \disponslide{4-}{A(r)}  & \disponslide{4-}{= \set{b, a}   } \\
                    \disponslide{5-}{B}     & \disponslide{5-}{= \bigcap \set{A(q), A(p), A(r)}} \\
                    \disponslide{5-}{}      & \disponslide{5-}{= \set{a, b}} \\
                    \disponslide{5-}{O_B}   & \disponslide{5-}{= \set{r, c}} \\
            \end{align*}
        \end{column}

        \begin{column}{0.4\textwidth}
            \includegraph{1}{fig9-1}
            \includegraph{2}{fig9-2}
            \includegraph{3}{fig9-3}
            \includegraph{4}{fig9-4}
            \includegraph{5-}{fig9-5}
        \end{column}
        \setbeamercovered{}
    \end{columns}
\end{frame}

\begin{frame}
    \begin{block}{Problem}
        BFS needs to start at left-most term $B_0$
    \end{block}

    \begin{block}{Solution}
        Process program depth-first, mark when processed
    \end{block}
\end{frame}

\begin{frame}
    \begin{itemize}
        \item If $v$ is fallthrough $\Rightarrow$ mark and continue
        \item If $v$ is conditional $\Rightarrow$ is $B_0$ and $B$ are marked
        \item If $v$ is marked $\Rightarrow$ continue
    \end{itemize}

    \begin{block}{Note}
        May lead to expressions being processed "out of order"
    \end{block}
\end{frame}

\begin{frame}[fragile]
    \begin{block}{When masking happens}
        \lstinline{* || true}

        \lstinline{* && false}
    \end{block}
\end{frame}

\begin{frame}[fragile]
    \begin{block}{Observation}
        Boolean expression are are isomorphic under the operator
    \end{block}

    \begin{columns}
        \begin{column}{0.4\textwidth}
            \centering
            \includegraphics[
                width   = \linewidth,
                height  = 0.8\textheight,
                keepaspectratio,
            ]{graph/fig10-1.pdf}

            \lstinline{a || b}
        \end{column}
        \begin{column}{0.4\textwidth}
            \centering
            \includegraphics[
                width   = \linewidth,
                height  = 0.8\textheight,
                keepaspectratio,
            ]{graph/fig10-2.pdf}

            \lstinline{a && b}
        \end{column}
    \end{columns}
\end{frame}

\begin{frame}
    \begin{block}{Proposition}
        Boolean expression are are isomorphic under the operator
    \end{block}

    \begin{block}{Proof}
        De Morgan's Laws
        \begin{align*}
            \neg(P \wedge Q) & \equiv \neg P \vee \neg Q \\
            \neg(P \vee Q) & \equiv \neg P \wedge \neg Q \\
        \end{align*}
    \end{block}

    \begin{block}{Implication}
        We don't need to know the operator, only the graph shape
    \end{block}
\end{frame}

\begin{frame}[fragile]
    \begin{block}{When masking happens (in CFG)}
        When a value $c$ is changed (taking a different edge at $v_c$) and
        we still end in the same outcome node
    \end{block}

    \lstinline{f = a || b}

    \begin{columns}
        \begin{column}{0.3\textwidth}
            \centering
            \includegraphics[
                width  = \linewidth,
                height = \textheight,
                keepaspectratio,
            ]{graph/fig11-1.pdf}
            \lstinline{f 1 0}
        \end{column}

        \begin{column}{0.3\textwidth}
            \centering
            \includegraphics[
                width  = \linewidth,
                height = \textheight,
                keepaspectratio,
            ]{graph/fig11-2.pdf}
            \lstinline{f 1 1}
        \end{column}

        \begin{column}{0.3\textwidth}
            \centering
            \includegraphics[
                width  = \linewidth,
                height = \textheight,
                keepaspectratio,
            ]{graph/fig11-3.pdf}
        \lstinline{f 0 1}
        \end{column}
    \end{columns}
\end{frame}

\begin{frame}
    \begin{block}{Observation}
        Masking happens at nodes with \emph{multiple} predecessors
    \end{block}

    \begin{block}{Implication}
        Multiple predecessors means short circuiting edge
    \end{block}

    \begin{block}{Implication}
        We know where to start searching
    \end{block}
\end{frame}

\begin{frame}
    \begin{block}{Association}
        \begin{align*}
            P \wedge (Q \wedge R) & \equiv (P \wedge Q) \wedge R \\
            P \vee (Q \vee R) & \equiv (P \vee Q) \vee R \\
        \end{align*}
    \end{block}

    \begin{block}{Implication}
        We can re-write expressions to an \emph{alternating} form

        \begin{align*}
            (A \vee B) \vee ((C \wedge D) \vee E) \\
            A \vee B \vee (C \wedge D) \vee E     \\
        \end{align*}
    \end{block}
\end{frame}

\begin{frame}
    \begin{block}{Observation}
        Masking propagates until the operator changes
    \end{block}

    \begin{block}{}
        \begin{equation*}
            A \wedge (B \vee C \vee D)
        \end{equation*}

        $D = t$ masks $B$, $C$, but not $A$
    \end{block}
\end{frame}

\begin{frame}
    \begin{block}{Subexpressions can mask}
        \begin{equation*}
            A \wedge (B \vee C)
        \end{equation*}
        $C = f$ masks $A$, but not $B$

        \begin{equation*}
            F = B \vee C
        \end{equation*}
        \begin{equation*}
            A \wedge F
        \end{equation*}
    \end{block}
\end{frame}

\begin{frame}
    \begin{block}{Observation}
        The \emph{last term} $B_\Omega$ in a subexpression can short circuit in
        the superexpression
    \end{block}
\end{frame}

\begin{frame}
    \begin{equation*}
            a \wedge b
        \wedge
            (c \vee d)
        \wedge
            e \wedge f
        \wedge
            (g
            \vee
                (h \wedge i)
            \vee
            j)
        \wedge
            k
    \end{equation*}

    \begin{tikzpicture}[node distance = 2em]
        \node (o1)               {$\wedge$};
        \node (o2) [below of=o1] {$\vee$};
        \node (o3) [below of=o2] {$\wedge$};

        \node (a)  [right of=o1] {$a$};
        \node (b)  [right of=a]  {$b$};
        \node (1)  [right of=b]  {};
        \node (c)  [below of=1]  {$c$};
        \node (d)  [right of=c]  {$d$};
        \node (2)  [above of=d]  {};
        \node (e)  [right of=2]  {$e$};
        \node (f)  [right of=e]  {$f$};
        \node (3)  [right of=f]  {};
        \node (g)  [below of=3]  {$g$};
        \node (4)  [right of=g]  {};
        \node (h)  [below of=4]  {$h$};
        \node (i)  [right of=h]  {$i$};
        \node (6)  [right of=i]  {};
        \node (j)  [above of=6]  {$j$};
        \node (7)  [right of=j]  {};
        \node (k)  [above of=7]  {$k$};
        \node (x)  [right of=k]  {};

        \only<2>{
            \draw [->] (c.{north}) to [bend left  = 30] (e.{west});
            \draw [->] (h.{north}) to [bend left  = 30] (j.{west});
        }

        \only<3>{
            CONTINUE HERE
            \draw [->] (d.{north}) to [bend left  = 30] (e.{west});
        }
    \end{tikzpicture}

    \only<2>{
        $c = true, h = true$
    }

    \only<3>{
        $d = false$
    }

\end{frame}

\begin{frame}
    \begin{equation*}
        Succ(B_\Omega) = O_B
    \end{equation*}

    \begin{block}{Observation}
        \begin{itemize}
            \item Either short-circuit right operands
            \item Evaluate next operand
        \end{itemize}
    \end{block}

    \begin{block}{Implication}
        If one edge is a short circuiting edge, the other must be a masking
        edge
    \end{block}
\end{frame}

\begin{frame}
    \begin{block}{Problem}
        Given a pair of incoming edges, which is masking and which is short
        circuting?
    \end{block}

    \begin{block}{Requirement}
        An ordering $v_n < v_m$ if $v_n$ is a left operand and $v_m$ is
        a right operand in the same expression
    \end{block}

    \begin{block}{Solution}<2>
        Topological sort
    \end{block}
\end{frame}

\begin{frame}
    Given $\set{v_n, v_m} = Preds(v), v_n < v_m$ then
    \begin{align*}
        v_n &= S_\Omega \\
        O_S &= Succ(v_n) \\
    \end{align*}
    where $S$ is a subexpression of $B$ ($S \subset B$)

    \begin{block}{Implication}
        When $(v_m, v)$ is taken, $S$ are masked
    \end{block}
\end{frame}

\begin{frame}
    \begin{columns}
        \begin{column}{0.5\textwidth}
            \begin{align*}
                { v_n, v_m } &= Preds(v) \\
                v_n &= S_\Omega \\
                O_S &=  Succ(v_n)
            \end{align*}
        \end{column}
    \end{columns}
\end{frame}

\end{document}
